\section{OS Foundation}
\label{sec:foundation}

Every operating system needs a foundation: a kernel that talks directly to
hardware, a startup system, a package manager, and a base set of system tools.
Building all of this from scratch would take decades and produce something
inferior to what already exists. This project therefore builds on top of an
existing Linux distribution that provides the entire foundation layer, and adds
everything above it: the desktop environment, event system, knowledge graph,
and AI layer. The distribution provides the load-bearing walls; the layout and
everything inside them is custom.

\subsection{Base Distribution}
\label{sec:base-distro}

The base distribution has not yet been finalized; two candidates are under
evaluation and the decision will be made before Phase~1 begins. Both use the
\emph{RPM} package format, which means the build pipeline and tooling are
compatible with either choice.

The requirements for the base are deliberately narrow. Packages need to be
current enough for modern development toolchains and recent hardware support,
but not so fresh that unvalidated changes arrive daily. Governance should not
be fully concentrated in a single commercial entity. \emph{Btrfs} must be
available as a first-class filesystem, since the update strategy described in
Section~\ref{sec:updates} depends on it. Driver and package coverage must be
broad enough to support consumer hardware without manual intervention.

\paragraph{Candidate A: Fedora.}
Fedora is a community distribution backed by Red Hat infrastructure but
governed independently through the Fedora Council, which makes binding
decisions on project direction without requiring Red Hat sign-off. Packages
track upstream closely and a new release ships every six months, making it
one of the more current stable distributions available. Btrfs has been the
default filesystem since Fedora~33, and the Flatpak and Wayland ecosystems
are better supported here than on most other distributions.

\begin{table}[htbp]
\caption{Fedora: evaluation summary}
\label{tab:fedora}
\begin{tabularx}{\linewidth}{XX}
\toprule
\textbf{Strengths} & \textbf{Weaknesses} \\
\midrule
Independent governance via Fedora Council &
  Six-month release cycle requires periodic major upgrades \\
Current packages, strong hardware support &
  Red Hat/IBM infrastructure in the background \\
Btrfs default, large ecosystem &
  Major upgrades need more integration testing than rolling updates \\
Strong Flatpak and Wayland support & \\
\bottomrule
\end{tabularx}
\end{table}

\paragraph{Candidate B: OpenSUSE Slowroll.}
Slowroll occupies the space between Tumbleweed, which is fully rolling and
ships new packages daily, and Leap, which follows a fixed annual release
cycle. It pulls packages from Tumbleweed after a validation delay of several
weeks, giving a cadence that is current enough for modern toolchains while
providing a buffer against regressions. Btrfs and Snapper are more deeply
integrated on OpenSUSE than anywhere else in the Linux ecosystem, which is
a genuine advantage for the update strategy.

\begin{table}[htbp]
\caption{OpenSUSE Slowroll: evaluation summary}
\label{tab:slowroll}
\begin{tabularx}{\linewidth}{XX}
\toprule
\textbf{Strengths} & \textbf{Weaknesses} \\
\midrule
Best Btrfs and Snapper integration in the ecosystem &
  SUSE owned by EQT Partners (private equity); governance risk \\
Rolling with validation delay; good update cadence &
  EQT ownership weakens European sovereignty narrative \\
Open Build Service is best-in-class for RPM packaging &
  Smaller community and hardware coverage than Fedora \\
European roots and presence & \\
\bottomrule
\end{tabularx}
\end{table}

The current preference is Fedora, primarily on governance grounds. The
EQT ownership of SUSE is not a technical problem, but it is a liability when
applying for European software independence funding through programs such as
the Sovereign Tech Fund or NGI: funders will scrutinize the ownership
structure and the connection to a US private equity firm sitting above a
nominally European project is a difficult position to defend. Fedora's
governance is cleaner for that argument. Both distributions use RPM packaging,
so the Open Build Service can target either without changes to the build
pipeline.

\subsection{Update Strategy}
\label{sec:updates}

Updates on most operating systems involve some degree of risk: something
breaks, the user does not know why, and getting back to a working state
requires either expertise or luck. This system addresses that with a
layered approach that makes the worst-case outcome a single keypress rather
than an afternoon in recovery mode.

The approach is a mutable filesystem with automatic rollback protection.
\emph{Mutable} means the filesystem works the way Linux filesystems always
have: software can be installed freely, system files can be modified, and
standard developer tools work without restrictions. This distinguishes it from
the immutable distributions that have gained popularity recently, where the
root filesystem is read-only and changes require going through a specific
toolchain. Immutable systems have genuine safety benefits, but they break
workflows that the target audience relies on, including installing language
runtimes, building projects from source, and running development servers
locally. The safety benefits of immutability are achievable through snapshot
tooling without those costs.

Figure~\ref{fig:update-flow} shows the full update flow.

\begin{figure}[htbp]
\centering
\begin{tikzpicture}[
  block/.style={
    draw, thick, rectangle, rounded corners=3pt,
    minimum width=5cm, minimum height=0.85cm,
    font=\small, align=center, inner sep=5pt
  },
  decision/.style={
    draw, thick, diamond, aspect=3,
    minimum width=5cm, minimum height=0.85cm,
    font=\small, align=center, inner sep=4pt
  },
  outcome/.style={
    draw, thick, rectangle, rounded corners=3pt,
    minimum width=3.8cm, minimum height=0.85cm,
    font=\small, align=center, inner sep=5pt
  },
  arr/.style={->, thick},
]

\def\sep{1.15}

\node[block]    (avail)    at (0, 4*\sep) {Update available};
\node[block]    (snap)     at (0, 3*\sep) {Pre-update snapshot (Snapper)};
\node[block]    (apply)    at (0, 2*\sep) {Update applied};
\node[block]    (boot)     at (0, \sep)   {First boot after update};
\node[decision] (check)    at (0, 0)      {Health check};
\node[outcome]  (ok)       at (-3.2, -2*\sep) {Continue normally};
\node[outcome]  (fail)     at ( 3.2, -2*\sep) {Auto-rollback + notify};

\draw[arr] (avail) -- (snap);
\draw[arr] (snap)  -- (apply);
\draw[arr] (apply) -- (boot);
\draw[arr] (boot)  -- (check);
\draw[arr] (check.west)  -| node[pos=0.2, above, font=\footnotesize] {pass} (ok);
\draw[arr] (check.east)  -| node[pos=0.2, above, font=\footnotesize] {fail} (fail);

\end{tikzpicture}
\caption{Update flow. A Snapper snapshot is created before every update.
If the post-update health check fails, the system rolls back to the
snapshot automatically.}
\label{fig:update-flow}
\end{figure}

\paragraph{Btrfs and Snapper.}
\emph{Btrfs} is a Linux filesystem with a built-in copy-on-write mechanism:
when a file changes, the old version is preserved rather than overwritten,
which makes filesystem-level snapshots nearly free in terms of disk space.
\emph{Snapper} is the snapshot management tool that sits on top of Btrfs and
handles creating, listing, and restoring snapshots automatically. The
filesystem is divided into subvolumes, with \texttt{/} and \texttt{/home}
on separate subvolumes so that a system rollback never touches personal data.

\paragraph{Pre-update snapshots.}
Before any update is applied, Snapper creates a snapshot of the current
system state. This happens automatically and requires no user action; it is a
precondition for the update, not an optional step.

\paragraph{Bootloader integration.}
If an update produces a system that fails to boot, the bootloader detects the
failure and presents the previous snapshot as a recovery option on the next
start. The user is not asked to open a terminal or enter recovery mode; they
see a clear prompt with a single keypress to continue on the last working
version.

\paragraph{Post-update health check.}
On the first boot after an update, a small background service verifies that
critical components are functional: the display system initializes correctly,
network is reachable, and no fatal errors appear in the service log. If the
check fails, the system rolls back to the pre-update snapshot and notifies the
user with a plain-language explanation. This catches the case where an update
boots but leaves the system in a degraded state that would not be obvious
until the user runs into something broken.

\paragraph{Snapshot retention.}
To prevent old snapshots from accumulating and consuming disk space,
Snapper removes snapshots older than seven days automatically while always
retaining the five most recent, regardless of age. Both values are adjustable
in Settings.

\subsection{System Infrastructure}
\label{sec:infra}

\emph{systemd} is used as-is with no replacements planned. It handles service
startup ordering, dependency management, and automatic restarts for crashed
services, and it is standard across both candidate distributions.
\texttt{systemd-homed} manages user home directories with built-in encryption,
which feeds into the account and profile system described in
Section~\ref{sec:platform}.

Package management at the base level uses the native manager for whichever
distribution is chosen: \texttt{dnf} for Fedora, \texttt{zypper} for
OpenSUSE. First-party project packages are built and distributed through the
\emph{Open Build Service} (OBS), a platform that compiles packages for
multiple distributions and CPU architectures from a single set of build
instructions and handles package signing and repository hosting. Store
applications are distributed as \emph{Flatpaks}, sandboxed packages that
include their own dependencies and run isolated from the rest of the system,
while system-level components ship as native RPMs. The Store infrastructure
handles both formats transparently.

\subsection{Hardware Requirements}
\label{sec:hardware}

Three hard requirements determine the minimum supported hardware. They are
not adjustable without significant changes to the security and graphics stack,
so they are treated as constraints rather than recommendations. In practice
they exclude hardware older than roughly 2018--2019.

\begin{table}[htbp]
\caption{Minimum hardware requirements}
\label{tab:hardware}
\begin{tabularx}{\linewidth}{lX}
\toprule
\textbf{Requirement} & \textbf{Reason} \\
\midrule
Linux kernel 6.6 or newer &
  Required for the eBPF features used by the event monitoring system and for
  current security capabilities. Standard on both candidate distributions. \\
Vulkan-capable GPU &
  Required for hardware-accelerated rendering and Windows compatibility.
  Covered automatically on AMD and Intel; Nvidia requires the proprietary
  driver, with installer guidance provided. \\
systemd 255 or newer &
  Required for full \texttt{systemd-homed} functionality used by the account
  system. \\
\bottomrule
\end{tabularx}
\end{table}

\subsection{Configuration Format}
\label{sec:config}

All system components and first-party applications use \emph{TOML} as the
configuration format. TOML is a human-readable format designed to be
unambiguous and easy to edit manually, without the indentation pitfalls of
YAML or the verbosity of JSON. Using it consistently across every component
means that a user or developer who understands one configuration file
understands all of them, and it means that tooling for reading, validating,
and migrating configuration only needs to be written once.

Personal settings live at \texttt{\textasciitilde/.config/<app-id>/config.toml}
and system-wide settings at \texttt{/etc/<service-name>/config.toml}. Some
third-party applications expect other formats; a small helper bridges this
without moving the source of truth:

\begin{lstlisting}[language=bash]
# Export as environment variables for a legacy app
eval $(os-config-helper export \
    --toml ~/.config/myapp/config.toml --format env)

# Export as JSON for tools that require it
os-config-helper export \
    --toml ~/.config/myapp/config.toml --format json
\end{lstlisting}

The TOML file is always the authoritative source; the helper is a read-only
translation layer with no business logic of its own.
